\documentclass[11pt]{article}

\usepackage[utf8]{inputenc}
\usepackage[T1]{fontenc}
\usepackage{lmodern}
\usepackage{geometry}
\usepackage{amsmath,amssymb,amsfonts,amsthm,mathtools}
\usepackage{bm}
\usepackage{enumitem}
\usepackage{microtype}
\usepackage{hyperref}
\usepackage{titlesec}
\usepackage{booktabs}
\usepackage{array}
\usepackage{setspace}
\usepackage{mathrsfs}
\usepackage{physics}

\geometry{margin=1in}
\hypersetup{colorlinks=true, linkcolor=black, urlcolor=blue, citecolor=black}
\setlist[enumerate]{topsep=0.4em,itemsep=0.35em}
\setlist[itemize]{topsep=0.3em,itemsep=0.25em}
\titleformat{\section}{\large\bfseries}{\thesection.}{0.5em}{}
\titleformat{\subsection}{\normalsize\bfseries}{\thesubsection.}{0.5em}{}
\setstretch{1.06}

\newcommand{\Aether}{\AE{}ther}
\newcommand{\Aflow}{\AE{}ther-flow}
\newcommand{\TheoryName}{\Aflow{} Interpretation of Relativity}
\newcommand{\TheoryProgram}{\Aether{} / \Aflow{} framework}

\newcommand{\CompletedMetric}{g^{\sharp}}
\newcommand{\MatterSpace}{\Psi}

\newcommand{\Car}{\mathrm{car}}
\newcommand{\Cur}{\mathrm{cur}}
\newcommand{\Hom}{\mathrm{hom}}

\newcommand{\ForSpace}[1]{\mathcal I_{#1}^{\mathrm{for}}}
\newcommand{\PrimShell}{\mathcal S_{\mathrm{prim}}}
\newcommand{\MetricOut}{\mathscr G_{\mathrm{out}}}
\newcommand{\MatterOut}{\mathscr T_{\mathrm{out}}}

\newcommand{\ProtoSectionBody}[1]{\mathcal B_{#1}^{\mathrm{psec}}}
\newcommand{\ProtoSectionShell}[1]{\mathcal Z_{#1}^{\mathrm{psec}}}
\newcommand{\PrePreProtoSectionVertical}[1]{V_{#1}^{\mathrm{pppresec}}}
\newcommand{\PrePreProtoSectionResidualBody}[1]{\mathcal D_{#1}^{\mathrm{pppresec}}}

\newcommand{\ReducedStructure}{\mathfrak S^{\mathrm{disp,resp,red}}}
\newcommand{\CurrentBodyImageCore}{\mathfrak S^{\mathrm{disp,resp,cbimg}}}
\newcommand{\ImageProjection}[1]{\Pi_{#1}^{\mathrm{disp,resp,cbimg}}}
\newcommand{\CanonicalLift}[1]{J_{#1}^{\mathrm{disp,resp,can}}}
\newcommand{\CanonicalImageBody}[1]{\mathcal U_{#1}^{\mathrm{disp,resp,cbimg}}}
\newcommand{\CanonicalImageResponseLaw}[1]{\mathscr Q_{#1}^{\mathrm{disp,resp,cbimg}}}
\newcommand{\CanonicalImageShell}[1]{\mathcal Z_{#1}^{\mathrm{disp,resp,cbimg}}}
\newcommand{\CoreForcedBodyResponse}[1]{\widehat{\mathscr R}_{#1}^{\mathrm{disp,resp,vis}}}
\newcommand{\VisibleResponseOperator}[1]{S_{#1}^{\mathrm{disp,resp,vis}}}

\newcommand{\CompatibleCarrier}[1]{\widetilde U_{#1}^{\mathrm{disp,resp}}}
\newcommand{\CompatibleProjection}[1]{\widetilde P_{#1}^{\mathrm{disp,resp}}}
\newcommand{\CompatibleOperator}[1]{\widetilde K_{#1}^{\mathrm{disp,resp}}}
\newcommand{\CompatibleKernel}[1]{\widetilde N_{#1}^{\mathrm{disp,resp}}}
\newcommand{\CompatibleSchur}[1]{\widetilde M_{#1}^{\mathrm{disp,resp}}}
\newcommand{\CompatibleLift}[1]{\widetilde J_{#1}^{\mathrm{disp,resp,can}}}
\newcommand{\CompatibleResponse}[1]{\widetilde{\mathscr R}_{#1}^{\mathrm{disp,resp,vis}}}
\newcommand{\MinimalReducedStructure}{\mathfrak S^{\mathrm{disp,resp,red,min}}}

\newtheorem{definition}{Definition}
\newtheorem{proposition}{Proposition}
\newtheorem{remark}{Remark}
\newtheorem{corollary}{Corollary}
\newtheorem{theorem}{Theorem}

\title{\textbf{The \TheoryName{}}\\[0.4em]
\large Primitive-Reservoir Observer-Reduction\\
\large Section-Centered Hidden-Response\\
\large Reduced Projection-Response Substrate Structure\\
\large Necessity / Uniqueness Theorem}
\author{}
\date{}

\begin{document}

\maketitle

\begin{abstract}
The current active-primary theorem already proves that, once a benchmark-faithful reduced projection-response substrate structure is fixed, the full canonical current-body image core is canonically derived and uniquely forced. The next honest move would ideally derive that reduced structure from still more primitive explicit substrate structure, but no genuinely cleaner deeper package is presently in hand without inventing another same-output presentation below the new frontier. The present note therefore takes the disciplined fallback explicitly allowed by the routing layer.

For each sector it proves that once the fixed full canonical current-body image core is held fixed, every benchmark-faithful reduced projection-response substrate structure compatible with that core induces the same visible quadratic response operator on the fixed section-centered displacement fibers and differs from the canonical minimal representative only by \(K\)-orthogonal inert kernel directions. Equivalently, the benchmark-facing content of the reduced structure is necessary and unique up to kernel-stable equivalence: carrier-level alternatives that preserve the same current image core carry no additional benchmark-facing freedom.

The benchmark boundary remains fixed throughout. The one operative metric is still exactly \(\CompletedMetric\), the ordinary matter route is unchanged, there is no second operative metric, the low-energy matter law is not enlarged, and no stronger promotion language is licensed. The positive result is therefore narrower than a new deeper derivation and stronger than another same-output relay: any honest deeper derivation of the fixed current image core must pass through the same reduced projection-response substrate structure, equivalently through its canonical minimal visible-response representative.
\end{abstract}

\section{Introduction}

The active derivational program of \TheoryName{} has now reached a point where the next benchmark-facing move is clear in principle but not yet clean in direct form. The current active-primary theorem already showed that the full canonical current-body image core \(\CurrentBodyImageCore\) is no longer the live primitive datum on the section-centered hidden-response line. Once the reduced projection-response substrate structure is fixed, the whole image core is already forced. The honest next move would therefore be to derive that reduced structure itself from still more primitive explicit substrate structure.

At present scope, however, a new deeper package that does that cleanly would require introducing another explicit substrate presentation whose benchmark-facing role is not yet disciplined enough to count as a genuine gain rather than another same-output relay. The routing layer therefore leaves one narrower but still benchmark-facing fallback: prove a sharper compatibility, necessity, or uniqueness theorem that makes the reduced-structure derivation unavoidable. The present note takes exactly that route.

The key point is simple. Even below the full-image-core forcing frontier, one could still worry that different reduced projection-response substrate structures might force the same image core through different carrier-level presentations. If so, a later deeper derivation might try to route around the reduced structure by choosing another compatible carrier package. The theorem below rules out that escape. Once the fixed current image core is held fixed, every compatible reduced projection-response substrate structure yields the same visible quadratic response on the fixed displacement fibers, and the only remaining freedom is the addition of \(K\)-orthogonal kernel directions that are invisible to the displacement projection and contribute only inert excess energy.

\paragraph{Framework and claim boundary.}
\TheoryProgram{} denotes the broader \Aether{} / \Aflow{} framework built on the claims that the \Aether{} is the underlying four-dimensional substrate of reality, the \Aflow{} is its intrinsic ordered motion, observed three-dimensional space is the local experiential slice of that deeper substrate, S-time is the experienced order of change arising from matter, light, and the \Aflow{}, and observed expansion is the three-dimensional appearance of deeper four-dimensional ordered motion. Within that framework, \TheoryName{} denotes the exact-closure theory in which the effective gravitational dynamics are adopted to be exactly Einsteinian with universal matter coupling. In the active sequence, \emph{adoption} means use of that established relativistic dynamics without claiming substrate derivation, while \emph{derivation} is reserved for a first-principles recovery from explicit substrate variables.

The present note stays strictly inside that boundary. It does not introduce a second operative metric, it does not enlarge the low-energy matter law, and it does not claim that the final substrate microphysics has been identified. Its narrower positive role is to show that, below the current full-image-core forcing frontier, the reduced projection-response substrate structure is itself unavoidable at benchmark-facing scope up to inert kernel extension.

\section{Fixed Inputs}

\begin{definition}[Recorded primitive continuation data]
\label{def:fixed}
Fix the following continuation data from the active primitive-reservoir observer-reduction line.
\begin{enumerate}
    \item For each sector label \(\sigma\in\{\Car,\Cur,\Hom\}\), a primitive forcing shell
    \[
    \ForSpace{\sigma}.
    \]
    \item The product shell
    \[
    \PrimShell
    \equiv
    \ForSpace{\Car}\times \ForSpace{\Cur}\times \ForSpace{\Hom}.
    \]
    \item The recorded metric and ordinary-matter outputs
    \[
    \MetricOut:\PrimShell\to\{\text{observer metrics}\},
    \qquad
    \MatterOut:\PrimShell\times\MatterSpace\to\{\text{observer stress tensors}\},
    \]
    satisfying
    \begin{equation}
    \MetricOut(\mathbf w)=\CompletedMetric,
    \qquad
    \MatterOut(\mathbf w,\psi)
    =
    T_{\mu\nu}^{\mathrm{matter}}[\CompletedMetric,\psi]
    \label{eq:fixedoutputs}
    \end{equation}
    for every \(\mathbf w\in\PrimShell\) and every matter configuration \(\psi\in\MatterSpace\).
\end{enumerate}
\end{definition}

\begin{definition}[Fixed current image-core frontier data]
\label{def:frontier}
Fix \(\sigma\in\{\Car,\Cur,\Hom\}\). The active line already fixes the following current image-core frontier data.
\begin{enumerate}
    \item The current proto-section benchmark base \(\ProtoSectionBody{\sigma}\) and exact proto-section shell \(\ProtoSectionShell{\sigma}\).
    \item The current section-centered displacement body
    \[
    \PrePreProtoSectionResidualBody{\sigma}
    \subset
    \ProtoSectionBody{\sigma}\times\PrePreProtoSectionVertical{\sigma}
    \]
    whose fiber
    \[
    \PrePreProtoSectionResidualBody{\sigma}(y)
    \equiv
    \left\{
    v\in\PrePreProtoSectionVertical{\sigma}:(y,v)\in\PrePreProtoSectionResidualBody{\sigma}
    \right\}
    \]
    is nonempty, compact, convex, contains \(0\), and has nonempty interior for every \(y\in\ProtoSectionBody{\sigma}\).
    \item The fixed canonical current-body image core
    \[
    \CurrentBodyImageCore_{\sigma}
    \equiv
    \Bigl(
    \CanonicalImageBody{\sigma},
    \ImageProjection{\sigma},
    \CanonicalLift{\sigma},
    \CanonicalImageResponseLaw{\sigma},
    \CanonicalImageShell{\sigma}
    \Bigr),
    \]
    where, for every \(y\in\ProtoSectionBody{\sigma}\),
    \[
    \ImageProjection{\sigma}(y):
    \CanonicalImageBody{\sigma}(y)\to\PrePreProtoSectionResidualBody{\sigma}(y)
    \]
    is a bijection with inverse \(\CanonicalLift{\sigma}(y)\), the response law \(\CanonicalImageResponseLaw{\sigma}(y,\cdot)\) is smooth and fiberwise quadratic, and the exact-shell image is
    \[
    \CanonicalImageShell{\sigma}
    =
    \left\{
    (y,0):y\in\ProtoSectionShell{\sigma}
    \right\}.
    \]
    \item Benchmark faithfulness, meaning preservation of \eqref{eq:fixedoutputs}, no second operative metric, no enlarged low-energy matter law, no change to the ordinary matter route, and no premature promotion from adoption to completed derivation.
\end{enumerate}
\end{definition}

\begin{proposition}[The fixed current image core determines a unique visible response operator]
\label{prop:visible}
For each \(y\in\ProtoSectionBody{\sigma}\), define
\begin{equation}
\CoreForcedBodyResponse{\sigma}(y,v)
\equiv
\CanonicalImageResponseLaw{\sigma}\bigl(y,\CanonicalLift{\sigma}(y)v\bigr)
\qquad
\text{for }
v\in\PrePreProtoSectionResidualBody{\sigma}(y).
\label{eq:fixedresponse}
\end{equation}
Then there exists a unique smooth family of positive self-adjoint operators
\[
\VisibleResponseOperator{\sigma}(y):
\PrePreProtoSectionVertical{\sigma}\to\PrePreProtoSectionVertical{\sigma}
\]
such that
\begin{equation}
\CoreForcedBodyResponse{\sigma}(y,v)
=
\frac{1}{2}
\ev{v,\VisibleResponseOperator{\sigma}(y)v}
\label{eq:visible}
\end{equation}
for every \((y,v)\in\PrePreProtoSectionResidualBody{\sigma}\).
\end{proposition}

\begin{proof}
For fixed \(y\), the function
\[
v\mapsto \CanonicalImageResponseLaw{\sigma}\bigl(y,\CanonicalLift{\sigma}(y)v\bigr)
\]
is quadratic on the fiber \(\PrePreProtoSectionResidualBody{\sigma}(y)\) because \(\CanonicalLift{\sigma}(y)\) is linear and the image response law is fiberwise quadratic. Since the fiber has nonempty interior and contains \(0\), that quadratic functional extends uniquely to a quadratic form on the whole vector space \(\PrePreProtoSectionVertical{\sigma}\). By polarization, there is a unique self-adjoint operator \(\VisibleResponseOperator{\sigma}(y)\) representing it as in \eqref{eq:visible}. Positivity follows because the fixed image-core response is nonnegative and vanishes only at \(v=0\), so the extended quadratic form is positive definite. Smoothness in \(y\) follows from the smooth dependence of \(\CanonicalLift{\sigma}\) and \(\CanonicalImageResponseLaw{\sigma}\) on \(y\).
\end{proof}

\begin{remark}
Proposition \ref{prop:visible} extracts exactly the datum that remains visible on the fixed section-centered displacement fibers once the full current image core is held fixed. The theorem below will show that every compatible reduced projection-response structure carries this same visible operator and no alternative benchmark-facing content beyond inert kernel directions.
\end{remark}

\section{Compatible Reduced Projection-Response Structures}

\begin{definition}[Compatible reduced projection-response substrate structure]
\label{def:compatible}
A \emph{benchmark-faithful compatible reduced projection-response substrate structure} for sector \(\sigma\) consists of the following data over the fixed frontier data of Definition \ref{def:frontier}.
\begin{enumerate}
    \item A finite-dimensional Euclidean compatible carrier
    \[
    \CompatibleCarrier{\sigma}.
    \]
    \item A smooth family of surjective linear displacement projections
    \[
    \CompatibleProjection{\sigma}(y):
    \CompatibleCarrier{\sigma}\to\PrePreProtoSectionVertical{\sigma},
    \qquad
    y\in\ProtoSectionBody{\sigma}.
    \]
    \item A smooth family of positive self-adjoint compatible response operators
    \[
    \CompatibleOperator{\sigma}(y):
    \CompatibleCarrier{\sigma}\to\CompatibleCarrier{\sigma},
    \qquad
    y\in\ProtoSectionBody{\sigma},
    \]
    satisfying the uniform coercive bound
    \begin{equation}
    \ev{u,\CompatibleOperator{\sigma}(y)u}
    \ge
    \mu_{\sigma}\,\|u\|^2
    \qquad
    \text{for all }
    y\in\ProtoSectionBody{\sigma},
    \ \ u\in\CompatibleCarrier{\sigma}
    \label{eq:coercive}
    \end{equation}
    for some constant \(\mu_{\sigma}>0\).
    \item The fixed exact-shell/current-body compatibility relation
    \[
    \left\{
    y\in\ProtoSectionBody{\sigma}:0\in\PrePreProtoSectionResidualBody{\sigma}(y)
    \right\}
    =
    \ProtoSectionShell{\sigma}.
    \]
    \item Benchmark faithfulness, meaning preservation of \eqref{eq:fixedoutputs}, no second operative metric, no enlarged low-energy matter law, no change to the ordinary matter route, and no premature promotion from adoption to completed derivation.
    \item Compatibility with the fixed current image-core frontier in the following precise sense. Define
    \begin{equation}
    \CompatibleSchur{\sigma}(y)
    \equiv
    \CompatibleProjection{\sigma}(y)
    \CompatibleOperator{\sigma}(y)^{-1}
    \CompatibleProjection{\sigma}(y)^*,
    \label{eq:schur}
    \end{equation}
    \begin{equation}
    \CompatibleLift{\sigma}(y)
    \equiv
    \CompatibleOperator{\sigma}(y)^{-1}
    \CompatibleProjection{\sigma}(y)^*
    \CompatibleSchur{\sigma}(y)^{-1},
    \label{eq:complift}
    \end{equation}
    and
    \begin{equation}
    \CompatibleResponse{\sigma}(y,v)
    \equiv
    \frac{1}{2}
    \ev{
    \CompatibleLift{\sigma}(y)v,
    \CompatibleOperator{\sigma}(y)\CompatibleLift{\sigma}(y)v
    }
    \label{eq:compresponse}
    \end{equation}
    for \((y,v)\in\PrePreProtoSectionResidualBody{\sigma}\). The compatibility requirement is
    \begin{equation}
    \CompatibleResponse{\sigma}(y,v)
    =
    \CoreForcedBodyResponse{\sigma}(y,v)
    \qquad
    \text{for every }
    (y,v)\in\PrePreProtoSectionResidualBody{\sigma}.
    \label{eq:compatibility}
    \end{equation}
\end{enumerate}
\end{definition}

\begin{remark}
Definition \ref{def:compatible} is intentionally the widest honest class below the current full-image-core forcing frontier. It does not assume literal equality with the currently recorded reduced structure. It assumes only the fixed displacement body, exact-shell discipline, and agreement of the induced visible response with the one already read out from the fixed current image core. The theorem below will show that this apparent freedom is benchmark-facingly illusory.
\end{remark}

\begin{proposition}[Compatible structures carry the same visible operator]
\label{prop:samevisible}
Let \(\widetilde{\ReducedStructure}_{\sigma}\) be a compatible reduced projection-response substrate structure in the sense of Definition \ref{def:compatible}. Then, for every \(y\in\ProtoSectionBody{\sigma}\),
\begin{equation}
\CompatibleResponse{\sigma}(y,v)
=
\frac{1}{2}
\ev{v,\CompatibleSchur{\sigma}(y)^{-1}v}
\qquad
\text{for all }
v\in\PrePreProtoSectionResidualBody{\sigma}(y),
\label{eq:schurresp}
\end{equation}
and
\begin{equation}
\CompatibleSchur{\sigma}(y)^{-1}
=
\VisibleResponseOperator{\sigma}(y).
\label{eq:samevisible}
\end{equation}
\end{proposition}

\begin{proof}
The operator \(\CompatibleSchur{\sigma}(y)\) is positive self-adjoint and invertible by the same surjectivity and coercivity argument used at the current full-image-core frontier: \(\CompatibleProjection{\sigma}(y)\) is surjective, so \(\CompatibleProjection{\sigma}(y)^*\) is injective, and \(\CompatibleOperator{\sigma}(y)^{-1}\) is positive definite. Using \eqref{eq:complift},
\begin{align*}
\CompatibleResponse{\sigma}(y,v)
&=
\frac{1}{2}
\ev{
\CompatibleOperator{\sigma}(y)^{-1}
\CompatibleProjection{\sigma}(y)^*
\CompatibleSchur{\sigma}(y)^{-1}v,
\CompatibleProjection{\sigma}(y)^*
\CompatibleSchur{\sigma}(y)^{-1}v
} \\
&=
\frac{1}{2}
\ev{
v,
\CompatibleSchur{\sigma}(y)^{-1}v
},
\end{align*}
which is \eqref{eq:schurresp}. By compatibility \eqref{eq:compatibility} and Proposition \ref{prop:visible},
\[
\frac{1}{2}
\ev{v,\CompatibleSchur{\sigma}(y)^{-1}v}
=
\frac{1}{2}
\ev{v,\VisibleResponseOperator{\sigma}(y)v}
\]
for every \(v\) in the fiber \(\PrePreProtoSectionResidualBody{\sigma}(y)\). Because that fiber has nonempty interior, the two quadratic forms agree on the whole vector space \(\PrePreProtoSectionVertical{\sigma}\). Hence their representing self-adjoint operators are equal, proving \eqref{eq:samevisible}.
\end{proof}

\begin{definition}[Canonical minimal reduced representative]
\label{def:minimal}
For sector \(\sigma\), define the \emph{canonical minimal reduced representative}
\[
\MinimalReducedStructure_{\sigma}
\equiv
\Bigl(
\PrePreProtoSectionVertical{\sigma},
\mathrm{id}_{\PrePreProtoSectionVertical{\sigma}},
\VisibleResponseOperator{\sigma}
\Bigr)
\]
over the fixed displacement body \(\PrePreProtoSectionResidualBody{\sigma}\), exact-shell relation \(\{y:0\in\PrePreProtoSectionResidualBody{\sigma}(y)\}=\ProtoSectionShell{\sigma}\), and the benchmark-faithfulness conditions of Definition \ref{def:frontier}(4).
\end{definition}

\begin{remark}
The minimal representative keeps only what remains benchmark-facingly visible below the current full-image-core forcing frontier: the fixed displacement body together with the fixed positive quadratic visible-response operator. All ambient carrier directions beyond that will be shown to be inert kernel extensions.
\end{remark}

\section{Kernel-Extension Decomposition}

\begin{proposition}[Compatible canonical lift and kernel decomposition]
\label{prop:decomp}
Let \(\widetilde{\ReducedStructure}_{\sigma}\) be a compatible reduced projection-response substrate structure. For every \(y\in\ProtoSectionBody{\sigma}\), let
\[
\CompatibleKernel{\sigma}(y)\equiv\ker\CompatibleProjection{\sigma}(y).
\]
Then:
\begin{enumerate}
    \item
    \[
    \CompatibleProjection{\sigma}(y)\circ\CompatibleLift{\sigma}(y)
    =
    \mathrm{id}_{\PrePreProtoSectionVertical{\sigma}};
    \]
    \item
    \[
    \ev{
    \CompatibleLift{\sigma}(y)v,
    \CompatibleOperator{\sigma}(y)k
    }
    =
    0
    \]
    for every \(v\in\PrePreProtoSectionVertical{\sigma}\) and every \(k\in\CompatibleKernel{\sigma}(y)\);
    \item every \(u\in\CompatibleCarrier{\sigma}\) admits a unique decomposition
    \[
    u
    =
    \CompatibleLift{\sigma}(y)\bigl(\CompatibleProjection{\sigma}(y)u\bigr)
    +
    k,
    \qquad
    k\in\CompatibleKernel{\sigma}(y);
    \]
    \item the energy splits as
    \begin{equation}
    \ev{u,\CompatibleOperator{\sigma}(y)u}
    =
    \ev{
    \CompatibleProjection{\sigma}(y)u,
    \VisibleResponseOperator{\sigma}(y)\CompatibleProjection{\sigma}(y)u
    }
    +
    \ev{k,\CompatibleOperator{\sigma}(y)k}
    \label{eq:split}
    \end{equation}
    for the unique decomposition in item (3).
\end{enumerate}
\end{proposition}

\begin{proof}
Item (1) follows directly from \eqref{eq:complift}. For item (2), using Proposition \ref{prop:samevisible},
\begin{align*}
\ev{
\CompatibleLift{\sigma}(y)v,
\CompatibleOperator{\sigma}(y)k
}
&=
\ev{
\CompatibleProjection{\sigma}(y)^*
\VisibleResponseOperator{\sigma}(y)v,
k
} \\
&=
\ev{
\VisibleResponseOperator{\sigma}(y)v,
\CompatibleProjection{\sigma}(y)k
}
=0.
\end{align*}
For item (3), set
\[
k
\equiv
u-\CompatibleLift{\sigma}(y)\bigl(\CompatibleProjection{\sigma}(y)u\bigr).
\]
Then item (1) gives \(\CompatibleProjection{\sigma}(y)k=0\), so \(k\in\CompatibleKernel{\sigma}(y)\). Uniqueness follows because the projection fixes the visible part.

For item (4), write \(u=\CompatibleLift{\sigma}(y)v+k\) with \(v=\CompatibleProjection{\sigma}(y)u\). By item (2),
\[
\ev{u,\CompatibleOperator{\sigma}(y)u}
=
\ev{
\CompatibleLift{\sigma}(y)v,
\CompatibleOperator{\sigma}(y)\CompatibleLift{\sigma}(y)v
}
+
\ev{k,\CompatibleOperator{\sigma}(y)k}.
\]
The first term is \(2\CompatibleResponse{\sigma}(y,v)\), which equals \(2\CoreForcedBodyResponse{\sigma}(y,v)\) by compatibility and therefore \(\ev{v,\VisibleResponseOperator{\sigma}(y)v}\) by Proposition \ref{prop:visible}. This proves \eqref{eq:split}.
\end{proof}

\section{Reduced-Structure Necessity / Uniqueness}

\begin{theorem}[Reduced projection-response substrate structure is necessary and unique up to inert kernel extension]
\label{thm:main}
Fix the primitive continuation data of Definition \ref{def:fixed} and the current image-core frontier data of Definition \ref{def:frontier}. Then the reduced projection-response substrate structure below the current full-image-core forcing frontier is benchmark-facingly necessary and unique up to inert kernel extension. More precisely:
\begin{enumerate}
    \item the canonical minimal representative \(\MinimalReducedStructure_{\sigma}\) of Definition \ref{def:minimal} is a benchmark-faithful reduced projection-response substrate structure, and its induced current image-core package is benchmark-faithfully equivalent to the fixed current image core \(\CurrentBodyImageCore_{\sigma}\);
    \item for every compatible reduced projection-response substrate structure \(\widetilde{\ReducedStructure}_{\sigma}\), there is a unique smooth fiberwise linear isomorphism
    \[
    \Theta_{\sigma}(y):
    \PrePreProtoSectionVertical{\sigma}\oplus\CompatibleKernel{\sigma}(y)
    \longrightarrow
    \CompatibleCarrier{\sigma}
    \]
    given by
    \[
    \Theta_{\sigma}(y)(v,k)
    \equiv
    \CompatibleLift{\sigma}(y)v+k,
    \]
    which satisfies
    \begin{equation}
    \CompatibleProjection{\sigma}(y)\Theta_{\sigma}(y)(v,k)=v
    \label{eq:theta_proj}
    \end{equation}
    and
    \begin{equation}
    \ev{
    \Theta_{\sigma}(y)(v,k),
    \CompatibleOperator{\sigma}(y)\Theta_{\sigma}(y)(v',k')
    }
    =
    \ev{v,\VisibleResponseOperator{\sigma}(y)v'}
    +
    \ev{k,\CompatibleOperator{\sigma}(y)k'}
    \label{eq:theta_energy}
    \end{equation}
    for all \(v,v'\in\PrePreProtoSectionVertical{\sigma}\) and \(k,k'\in\CompatibleKernel{\sigma}(y)\);
    \item consequently, every compatible reduced projection-response substrate structure differs from the canonical minimal representative only by adjoining \(K\)-orthogonal kernel directions that are invisible to the displacement projection and inert at benchmark-facing scope;
    \item in particular, any deeper benchmark-faithful derivation that forces the fixed current full image core must derive a compatible reduced projection-response substrate structure and therefore the same canonical minimal representative. Another manuscript that changes only compatible kernel extensions does not change benchmark-facing status.
\end{enumerate}
\end{theorem}

\begin{proof}
For item (1), the carrier \(\PrePreProtoSectionVertical{\sigma}\) is finite dimensional, the identity map is surjective, and Proposition \ref{prop:visible} gives a positive self-adjoint operator \(\VisibleResponseOperator{\sigma}(y)\) on it. The fixed exact-shell relation and benchmark-faithfulness conditions are inherited directly from Definition \ref{def:frontier}. The induced canonical lift of the minimal representative is the identity map on \(\PrePreProtoSectionVertical{\sigma}\), so its image body is precisely \(\PrePreProtoSectionResidualBody{\sigma}(y)\). The resulting response law is \(\CoreForcedBodyResponse{\sigma}(y,v)\) by \eqref{eq:visible}. Therefore the map
\[
v\longmapsto \CanonicalLift{\sigma}(y)v
\]
from the minimal image body \(\PrePreProtoSectionResidualBody{\sigma}(y)\) to the fixed image body \(\CanonicalImageBody{\sigma}(y)\) intertwines the image-body bijection, the response law, and the exact-shell image. Hence the induced image-core package is benchmark-faithfully equivalent to the fixed current image core.

For item (2), Proposition \ref{prop:decomp}(3) already shows that every element of \(\CompatibleCarrier{\sigma}\) has a unique expression \(\CompatibleLift{\sigma}(y)v+k\), so \(\Theta_{\sigma}(y)\) is a linear isomorphism. Smoothness follows from the smoothness of \(\CompatibleLift{\sigma}\). Equation \eqref{eq:theta_proj} is immediate from Proposition \ref{prop:decomp}(1). Equation \eqref{eq:theta_energy} follows from Proposition \ref{prop:decomp}(2) and the energy-splitting identity \eqref{eq:split}. Uniqueness holds because the visible part is fixed by the projection and the kernel part is then fixed by subtraction.

Item (3) is the interpretation of item (2): the visible benchmark-facing datum is always the same operator \(\VisibleResponseOperator{\sigma}\) on the same displacement space, while the only remaining freedom lies in kernel directions annihilated by the displacement projection and separated orthogonally from the visible part by the compatible response operator.

Item (4) is immediate. Any deeper benchmark-faithful derivation of the fixed current image core must, at the point where it yields a reduced projection-response structure, produce a compatible member of Definition \ref{def:compatible}; items (1)--(3) then force the same canonical minimal representative. A manuscript that only changes compatible kernel extensions therefore changes no benchmark-facing datum below the current full-image-core forcing frontier.
\end{proof}

\begin{corollary}[Unavoidability of reduced-structure derivation]
\label{cor:unavoidable}
Below the current full-image-core forcing frontier, the reduced projection-response substrate structure is the unavoidable benchmark-facing bridge object: a deeper derivation can no longer honestly route around it through another compatible carrier presentation.
\end{corollary}

\begin{proof}
This is exactly item (4) of Theorem \ref{thm:main}.
\end{proof}

\section{Conclusion}

The positive result of this note is narrower than a new deeper derivation and stronger than another same-output relay. The current full-image-core forcing theorem already showed that the full canonical current-body image core \(\CurrentBodyImageCore\) is canonically derived once the reduced projection-response substrate structure is fixed. The present theorem now sharpens that gain one honest level further: once the fixed current image core is held fixed, every compatible reduced projection-response substrate structure carries the same visible quadratic response operator on the fixed displacement fibers and differs from the canonical minimal representative only by inert kernel directions.

The scope remains disciplined. The benchmark package is unchanged: the one operative metric is still exactly \(\CompletedMetric\), the ordinary matter route is unchanged, there is no second operative metric, and the low-energy matter law is not enlarged. Nothing here upgrades adoption to derivation or licenses stronger promotion language. The result is a benchmark-facing uniqueness theorem, not a claim of completed substrate microphysics.

The next open burden is therefore clear. The next benchmark-facing manuscript should derive or justify the reduced projection-response substrate structure itself from still more primitive explicit substrate structure, equivalently derive its canonical minimal visible-response representative, while preserving the same benchmark one-metric package, the same ordinary matter route, and the same claim boundary. Another same-output relay that merely adjoins inert kernel directions, another restoration of the full substrate-body presentation, or another replay of the full projection-operator core or the full displacement-response substrate law is not the clean next manuscript target at current scope.

\end{document}
